% ch04-poisson --- Poisson models: the baselines that must be beaten

\section{Why the ladder starts here}\label{ch04:sec:role}

The timing target T2 of \cref{ch:question} asks for the ground intensity
$\lam_s(t)$: how many deals, in which sector, next week, with calibrated
uncertainty. This chapter builds the two simplest answers --- the
homogeneous Poisson process, which asserts that nothing is valuable to know
beyond the level of activity, and the inhomogeneous Poisson process (NHPP)
with covariates, which asserts that the \emph{market clock} is valuable but
the market's own history is not. Their role in the design is adversarial.
Every richer timing model in this document --- the
history-modulated GLM previewed in \cref{ch04:sec:history}, the exponential
Hawkes machinery of \cref{ch:hawkes}, the interval-censored variants of
\cref{ch:censoring} --- is admitted into production only by beating a
well-covariated NHPP on held-out likelihood under the protocols of
\cref{ch:gof}. Our working hypothesis is that covariates carry most of the
held-out gain over a seasonal baseline, and that everything added on top
of them --- lags, excitation, cross-sector terms --- contributes far less;
the decomposition benchmarks planned in \cref{ch:results-synth} will test
this. A Hawkes fit whose covariate baseline
is impoverished will absorb the macro cycle into its kernel and report
spurious contagion; the only defense is a baseline that already knows the
macro cycle, which is what this chapter calibrates.

The second reason to start here is that this is the one rung whose
likelihood is \emph{exactly right} under the observation scheme of
\cref{ch:data}. At sector resolution, right censoring costs nothing
(\cref{ch04:sec:censoring} proves this), interval censoring by weekly
aggregation costs nothing (\cref{thm:ch04:glmequiv}), and sectors neither
enter late nor die, so (C2) and (C3) do not bind. Calibration discipline is
cheap at this rung. Every later deviation from it must be priced against a
baseline whose own accounting is exact.

\section{The two models and their assumptions}\label{ch04:sec:models}

Throughout, units $e$ are sectors $s \in \sectors$ (the primary case) or,
for the event-time path, tracked firms $i \in \firms$; $N_e$ is the
counting process of \cref{def:ch02:intensity} with at-risk indicator
$\atrisk_e$, and all likelihoods specialize the Jacod likelihood
\cref{eq:ch02:jacod} whose martingale estimation theory is developed in
\cref{ch:foundations}. The homogeneous model sets $\lam_e(t) \equiv \mub_e$:
conditionally on nothing, events arrive at a constant rate, the MLE is
$\hat\mub_e = N_e(\Tend) / \int_0^{\Tend}\atrisk_e$, and the model's only
message is the level. It is the null model for timing; its
calendar-dummy refinement (a rate per quarter or per season) is the
``seasonal baseline'' against which covariate gains are measured.
The inhomogeneous model is the first calibrated one, and we now fix its
assumptions as numbered contracts.

\begin{assumption}[Conditional Poisson structure]\label{asm:ch04:condpoisson}
Given the covariate paths, the intensity of $N_e$ is a deterministic
function of time: $\lam_e(t;\thetamodel) = \lam_e(t)$ measurable with
respect to the covariate history alone, with no dependence on the internal
history $\Hist_{t^-}$ of the deal process. Equivalently, conditionally on
covariates, $N_e$ is an inhomogeneous Poisson process: counts over disjoint
sets are independent, and the count over a set $A$ is Poisson with mean
$\int_A \atrisk_e(u)\lam_e(u)\dd u$ \citep{kingman1993,daley2003}.
\end{assumption}

\begin{assumption}[Log-linear intensity]\label{asm:ch04:loglin}
$\lam_s(t) = \exp\!\big(\bcoef\T \xcov_s(t)\big)$, where $\xcov_s(t) \in
\R^{d}$ stacks a sector intercept (equivalently, sector dummies), the
published as-of macro values of \cref{prot:ch02:pit}(a), and deterministic
seasonality and calendar terms; per-sector slopes are interactions of the
macro block with the sector dummies. The exponential link is the canonical
one; \cref{ch04:sec:link} defends the choice against its alternatives.
\end{assumption}

\begin{assumption}[Piecewise-constant weekly covariates]\label{asm:ch04:pwconst}
$\xcov_s(\cdot)$ is constant on the cells of a partition of $[0,\Tend]$
into weekly bins: macro as-of series are step functions by construction
(they jump at publication dates, which we align to the weekly grid),
calendar terms are constant within a week, and sector dummies are constant
always. Consequently $\lam_s(\cdot)$ is a step function.
\end{assumption}

\begin{assumption}[External covariates]\label{asm:ch04:external}
The covariates of \cref{asm:ch04:loglin} are \emph{external} in the sense
of \citet{kalbfleisch2002}: their evolution does not depend on the deal
stream. Macro series, calendar terms, and sector dummies qualify. Lagged
deal counts, exponentially filtered event streams, and any functional of
$\Hist_{t^-}$ do \emph{not}; adding them changes the model class
(\cref{ch04:sec:history}) and voids \cref{asm:ch04:condpoisson}.
\end{assumption}

\begin{assumption}[Spanning]\label{asm:ch04:span}
Writing $c$ for the sector-week cells of \cref{asm:ch04:pwconst}, $w_c =
\int_{\text{cell}} \atrisk(u) \dd u$ for the exposure of cell $c$ and
$\xcov_c$ for its covariate vector, the family $\{\xcov_c : w_c > 0\}$
spans $\R^{d}$; equivalently $\sum_c w_c \xcov_c \xcov_c\T \succ 0$.
\end{assumption}

\Cref{asm:ch04:condpoisson} is a known fiction --- \cref{ch04:sec:overdisp}
measures its failure and upgrades the noise model --- while
\cref{asm:ch04:pwconst,asm:ch04:external} are believed exactly true for
the design of \cref{asm:ch04:loglin}, and \cref{asm:ch04:span} is checkable
from the design matrix before any fit.

\subsection{Why the exponential link}\label{ch04:sec:link}

Three properties recommend $\exp$. \emph{Positivity for free}: any $\bcoef \in
\R^{d}$ yields a valid intensity, so estimation is unconstrained smooth
optimization --- no boundary, no barrier, no projected gradient.
\emph{Concavity}: the log-likelihood is globally concave in $\bcoef$
(\cref{thm:ch04:concavity}), so the fitted object is a reproducible global
optimum rather than a local one an optimizer happened to find.
\emph{Multiplicative interpretation}: a unit move in covariate $j$
multiplies the rate by $e^{\bcoef_j}$ uniformly across sectors and weeks,
so macro effects read as rate ratios --- the natural language for
``tightening moved weekly late-stage deal flow down 12\%'' --- and sector
dummies give sector-proportional baselines. The log link is also
\emph{aggregation-coherent} under \cref{asm:ch04:pwconst}: binning to weeks
leaves the functional form intact with the bin width absorbed into the
intercept (\cref{thm:ch04:glmequiv}), which is precisely what makes the
weekly GLM an exact NHPP calibration rather than a discretization of one.

The alternatives are worth stating because the repository uses one of them.
The \emph{identity link} $\lam = \bcoef\T\xcov$ makes effects additive ---
deals per week per unit of covariate --- which is the natural scale for
excitation terms with a branching interpretation, and it tolerates
within-bin covariate variation exactly (the integral of a linear function
is linear in the averaged covariate, so no Jensen gap ever opens). Its
price is positivity: the likelihood is concave only on the polyhedron
$\{\bcoef : \bcoef\T\xcov_c > 0\ \forall c\}$, the MLE can sit on the
boundary where standard asymptotics fail \citep{selfliang1987}, and every
fit needs constraints. The \emph{softplus link} $\lam = \log(1 +
e^{\eta})$ keeps positivity free and grows only linearly for large $\eta$,
which reduces the exponential extrapolation risk of the log link outside
the training covariate support; it remains log-concave, but the
coefficients lose their rate-ratio reading and the fit leaves standard GLM
software. The stabilized sector layer of \cref{ch:hawkes} is a deliberate
hybrid: exponential link for the covariate baseline (this chapter's
object, with all three advantages), identity link for the added excitation
term (where the branching reading and the stability constraint live).

\begin{decision}{D4.1 --- exponential link for the covariate baseline}
The covariate-driven baseline intensity of every ground-process model in
this document uses the log-linear form of \cref{asm:ch04:loglin}.
Rationale: unconstrained concave estimation (\cref{thm:ch04:concavity}),
exact weekly-GLM equivalence (\cref{thm:ch04:glmequiv}), multiplicative
macro interpretation, and sign freedom for later history terms (the
exp-link route is one of the three legal ways to express inhibition;
\cref{ch:hawkes}). Reversal trigger: if held-out evaluation shows
exponential extrapolation failures --- predictions pinned at the
support-envelope guard of \modrepo{sector\_stability.py} on more than
isolated weeks --- the baseline link is switched to softplus and this
chapter's concavity and equivalence results are re-derived for it (both
survive in weakened form; the GLM packaging does not).
\end{decision}

\section{Likelihood, score, and the geometry of the optimum}\label{ch04:sec:likelihood}

Under \cref{asm:ch04:condpoisson,asm:ch04:loglin}, the Jacod likelihood
\cref{eq:ch02:jacod} specializes, for the pooled sector process on
$[0,\Tend]$, to
\begin{equation}\label{eq:ch04:loglik}
\loglik(\bcoef)
 \;=\;
 \sum_{m} \bcoef\T \xcov_{s_m}(t_m)
 \;-\;
 \sum_{s\in\sectors} \int_{0}^{\Tend} \atrisk_s(u)\,
 e^{\bcoef\T \xcov_s(u)} \dd u ,
\end{equation}
the sum running over observed deals $m$ with sector $s_m$ and time $t_m$.
Under \cref{asm:ch04:pwconst} the integral collapses over cells: with
$y_c$ the event count, $w_c$ the exposure, and $\xcov_c$ the covariate of
cell $c$,
\begin{equation}\label{eq:ch04:cellloglik}
\loglik(\bcoef)
 \;=\;
 \sum_{c} \Big( y_c\, \bcoef\T \xcov_c \;-\; w_c\, e^{\bcoef\T \xcov_c}
 \Big).
\end{equation}
Differentiating \cref{eq:ch04:cellloglik} (or \cref{eq:ch04:loglik}; the
cell form is the general one with the integral discretized exactly):
\begin{align}
\nabla \loglik(\bcoef)
 &\;=\;
 \sum_{c} \big( y_c - w_c\, e^{\bcoef\T \xcov_c} \big)\, \xcov_c ,
 \label{eq:ch04:score}\\
\nabla^2 \loglik(\bcoef)
 &\;=\;
 -\sum_{c} w_c\, e^{\bcoef\T \xcov_c}\, \xcov_c \xcov_c\T .
 \label{eq:ch04:hessian}
\end{align}
The score \cref{eq:ch04:score} is the familiar ``observed minus expected''
weighted by covariates; evaluated at the truth it is the terminal value of
the martingale $\int \xcov\, (\dd N - \atrisk\,\lam\,\dd t)$, which is the
engine of the asymptotics below. The Hessian \cref{eq:ch04:hessian} does
not involve the counts $y_c$ at all --- a canonical-link property: observed
information equals conditional expected information, so Newton's method,
Fisher scoring, and iteratively reweighted least squares are the same
algorithm here.

\begin{theorem}[Strict log-concavity under the exponential link]\label{thm:ch04:concavity}
Under \cref{asm:ch04:condpoisson,asm:ch04:loglin,asm:ch04:pwconst}, the
map $\bcoef \mapsto \loglik(\bcoef)$ of \cref{eq:ch04:cellloglik} is
concave on all of $\R^{d}$:
$\nabla^2\loglik(\bcoef) \preceq 0$ for every $\bcoef$. If in addition
\cref{asm:ch04:span} holds, it is strictly concave:
$\nabla^2\loglik(\bcoef) \prec 0$ for every $\bcoef$.
\end{theorem}

\begin{proof}
For any $v \in \R^{d}$,
\[
v\T \nabla^2 \loglik(\bcoef)\, v
 \;=\;
 -\sum_{c} w_c\, e^{\bcoef\T\xcov_c}\, \big(v\T \xcov_c\big)^2
 \;\le\; 0,
\]
since every summand is a product of nonnegative factors: concavity, with
no condition. Equality holds iff $v\T\xcov_c = 0$ for every cell with
$w_c > 0$ (the weights $w_c e^{\bcoef\T\xcov_c}$ are strictly positive on
exactly those cells). Under \cref{asm:ch04:span} no $v \ne 0$ annihilates
all exposed $\xcov_c$, so the quadratic form is strictly negative for all
$v \neq 0$ and all $\bcoef$. When \cref{asm:ch04:span} fails --- say two
macro series are exactly collinear over the training window --- the
likelihood is flat along the offending direction and $\bcoef$ is
identified only up to it; the fitted \emph{intensity} is still unique,
which is why ridge regularization plus inference on well-identified linear
combinations, rather than coefficient surgery, is the prescribed response
to the near-failure that smooth, correlated macro series actually produce.
\end{proof}

\begin{corollary}[Global optimum and solver guarantees]\label{cor:ch04:global}
Under the assumptions of \cref{thm:ch04:concavity} with
\cref{asm:ch04:span}, the MLE, when it exists, is the unique global
maximizer of $\loglik$, every stationary point is that maximizer, and any
ascent method with standard line-search or trust-region safeguards ---
Newton, Fisher scoring, L-BFGS --- converges to it from any starting
point. Two fits of the same data that report different optima are
evidence of a software defect, not of multimodality.
\end{corollary}

\begin{proof}
Strict concavity on $\R^d$ makes the superlevel sets convex and any
stationary point the unique global maximum; convergence of safeguarded
ascent methods on smooth concave objectives with nonsingular Hessian is
standard. The content worth stating is the last sentence: at this rung,
optimizer nondeterminism is a bug detector. This guarantee is exactly what
is forfeited --- first partially (block concavity), then wholly --- as
excitation parameters with free decays enter in \cref{ch:hawkes}.
\end{proof}

\paragraph{Asymptotics.}
At the true $\bcoef_0$ the score is a martingale integral with predictable
variation $\Fisher(\bcoef_0) = \sum_c w_c e^{\bcoef_0\T\xcov_c} \xcov_c
\xcov_c\T$ --- the Fisher information, equal to minus the Hessian
\cref{eq:ch04:hessian}. Rebolledo's martingale CLT then gives
$\hat\bcoef \approx \mathcal{N}\big(\bcoef_0, \Fisher^{-1}\big)$ as
information accumulates, through a longer window or more sectors; this is
the counting-process estimation theory of \citet{abgk1993} and
\citet{andersengill1982}, and the stationary-point-process version of
\citet{ogata1978}, specialized to the concave case. What matters
practically is that information is \emph{pooled}: a $d$-dimensional
$\bcoef$ needs adequate total information $\sum_c w_c
e^{\bcoef\T\xcov_c}\xcov_c\xcov_c\T$, not adequate counts in every
sector-week cell --- sparse cells are a separation risk
(\cref{prop:ch04:separation}), not an asymptotics failure per se.

\subsection{Existence: separation and sparse cells}\label{ch04:sec:separation}

Concavity guarantees uniqueness, not existence. The Poisson analogue of
logistic separation is the one genuine pathology of this rung.

\begin{proposition}[Nonexistence of the MLE under separation]\label{prop:ch04:separation}
Under the assumptions of \cref{thm:ch04:concavity} with
\cref{asm:ch04:span}, the MLE fails to exist iff there is a direction
$v \neq 0$ with
\begin{enumerate}[label=(\roman*)]
\item $v\T \xcov_c \le 0$ for every cell with $w_c > 0$, and
\item $v\T \xcov_c = 0$ for every cell with $y_c > 0$.
\end{enumerate}
Along any such $v$ the likelihood increases strictly and monotonically to
a finite supremum that is not attained; the fitted intensity on the cells
with $v\T\xcov_c < 0$ degenerates to zero.
\end{proposition}

\begin{proof}
Suppose such $v$ exists. By \cref{asm:ch04:span}, $v\T\xcov_c \ne 0$ for
at least one exposed cell, which by (ii) has $y_c = 0$ and by (i) has
$v\T\xcov_c < 0$. Along $\bcoef + \tau v$, the event term of
\cref{eq:ch04:cellloglik} is constant in $\tau$ by (ii), while each
compensator term $-w_c e^{\bcoef\T\xcov_c + \tau v\T\xcov_c}$ is
nondecreasing in $\tau$ by (i) and strictly increasing on the cell just
exhibited. Hence $\loglik$ increases strictly to the finite limit obtained
by deleting the vanishing compensator terms, and no finite $\bcoef$
attains it. Conversely, if no such $v$ exists, then along every ray
$\bcoef + \tau v$ either some exposed cell has $v\T\xcov_c > 0$, in which
case the compensator term drives $\loglik \to -\infty$ exponentially
(dominating the linearly growing event term), or all exposed cells have
$v\T\xcov_c \le 0$ with strict inequality at some event cell, in which
case the event term drives $\loglik \to -\infty$ linearly. So $\loglik$ is
coercive on $\R^d$ and, being continuous, attains its maximum.
\end{proof}

\begin{warning}[Sparse cells produce infinite coefficients]\label{warn:ch04:separation}
The separating directions of \cref{prop:ch04:separation} are not exotic:
any indicator covariate whose active cells contain \emph{zero events} ---
a small sector crossed with a rare calendar dummy, an interaction active
only in a few quiet weeks --- yields $v$ = minus that coordinate, and the
fitted coefficient runs to $-\infty$ (the MLE of a zero rate). At weekly
resolution with twelve sectors and rich calendar interactions this
\emph{will} occur. Remedies, in the order the repository applies them: the
ridge penalty already standard in the fitters (a strictly concave
perturbation restores existence unconditionally; the penalty must be
reported next to the fit); the Firth adjustment, maximizing $\loglik(\bcoef)
+ \tfrac{1}{2}\log\det\Fisher(\bcoef)$, which removes the leading-order
bias of the MLE and keeps estimates finite \citep{firth1993}, with the
rare-events analysis of \citet{kingzeng2001} as the companion warning
that unpenalized rates estimated from a handful of events are biased as
well as fragile; and design hygiene --- collapse or drop cells/dummies
that the training window cannot inform. Related but distinct is
\emph{extrapolation} risk: a log-linear intensity evaluated outside the
training covariate support is an exponential extrapolation even when the
MLE exists; the support-envelope winsorization in
\modrepo{sector\_stability.py} is the implemented guard, and train/test
covariate ranges are reported next to every held-out NLL
(\cref{ch:gof}).
\end{warning}

\section{The binned-GLM equivalence theorem}\label{ch04:sec:glm}

The repository fits weekly sector counts with a Poisson GLM. The following
theorem says this \emph{is} NHPP maximum likelihood --- exactly, not
approximately --- and identifies precisely what binning discards
(nothing, under this chapter's assumptions).

\begin{theorem}[Piecewise-constant NHPP $=$ Poisson GLM with log-exposure offset]\label{thm:ch04:glmequiv}
Let \cref{asm:ch04:condpoisson,asm:ch04:loglin,asm:ch04:pwconst} hold, so
that $\lam$ is constant, equal to $\lam_c = e^{\bcoef\T\xcov_c}$, on each
cell $c$ of the weekly partition, and let $w_c = \int_{\text{cell}}
\atrisk(u)\dd u$ be the (possibly partial) exposure of cell $c$. Then:
\begin{enumerate}[label=(\roman*)]
\item \emph{(Distribution.)} Conditionally on covariates and exposures,
the cell counts are independent with
$y_c \sim \mathrm{Poisson}\big(w_c\, e^{\bcoef\T \xcov_c}\big)$.
\item \emph{(Likelihood identity.)} The point-process log-likelihood
\cref{eq:ch04:loglik} and the log-likelihood of the Poisson GLM with log
link, design vectors $\xcov_c$, and offset $\log w_c$ differ by a constant
free of $\bcoef$:
\begin{equation}\label{eq:ch04:glm}
\loglik_{\mathrm{pp}}(\bcoef)
 \;=\;
 \sum_c \Big( y_c \log\!\big(w_c e^{\bcoef\T\xcov_c}\big)
 - w_c e^{\bcoef\T\xcov_c} - \log y_c! \Big)
 \;+\; \text{const}.
\end{equation}
\item \emph{(Sufficiency.)} The vector of cell counts $(y_c)_c$ is
sufficient for $\bcoef$: given the counts, the event times within each
cell are distributed as i.i.d.\ uniform draws on the at-risk subset of the
cell, a law free of $\bcoef$.
\end{enumerate}
Consequently the NHPP MLE, score, observed information, and all
likelihood-based inference coincide with those of the binned Poisson GLM;
fitting the weekly sector count GLM is an NHPP calibration, not a
discretization of one.
\end{theorem}

\begin{proof}
(i) is the independent-increments property of
\cref{asm:ch04:condpoisson} applied to the disjoint at-risk subsets of the
cells: the count over the at-risk subset $A_c$ of cell $c$ is Poisson with
mean $\int_{A_c} \lam_c \dd u = w_c \lam_c$
\citep{kingman1993}.

(ii) Because $\xcov$ is constant on cells, the event term of
\cref{eq:ch04:loglik} groups as $\sum_m \bcoef\T\xcov(t_m) = \sum_c y_c\,
\bcoef\T\xcov_c$, and the compensator integral collapses to $\sum_c w_c
e^{\bcoef\T\xcov_c}$, giving \cref{eq:ch04:cellloglik}. Expanding the GLM
log-likelihood on the right of \cref{eq:ch04:glm}:
\[
\sum_c \Big( y_c \log w_c + y_c\,\bcoef\T\xcov_c
 - w_c e^{\bcoef\T\xcov_c} - \log y_c! \Big)
 \;=\;
\loglik_{\mathrm{pp}}(\bcoef)
 + \sum_c \big( y_c \log w_c - \log y_c! \big),
\]
and the bracketed sum is a statistic of the data only.

(iii) Conditionally on $y_c = n$, the $n$ points of a Poisson process with
\emph{constant} rate on $A_c$ are distributed as the order statistics of
$n$ i.i.d.\ uniforms on $A_c$, for every value of the rate
\citep{kingman1993}; the conditional law is therefore free of $\bcoef$,
and sufficiency follows from the factorization
$\lik(\bcoef; \text{times}) = \lik(\bcoef; \text{counts}) \times
p(\text{times} \st \text{counts})$. The final claims follow by
differentiating an identity that holds up to an additive constant.
\end{proof}

\begin{corollary}[The repository's weekly GLM is an NHPP fit]\label{cor:ch04:repoglm}
With equal-width weekly bins and fully exposed sectors ($w_c \equiv w$),
the offset $\log w$ is absorbed into the sector intercepts, and the weekly
sector Poisson GLM --- \code{fit\_sector\_count\_model} of
\modrepo{sector\_stability.py} with \code{n\_lags=0}, and its vectorized
synthetic-world twin \code{fit\_sector\_glm\_fast(n\_lags=0)} --- computes
exactly the NHPP MLE of
\cref{asm:ch04:loglin,asm:ch04:pwconst}. Partial exposure (a unit entering
or leaving observation mid-week) is handled exactly by the offset, with no
change to the estimator's form.
\end{corollary}

Two remarks bound the theorem's reach. First, parts (i) and (iii) never
used the log link: the reduction of a piecewise-constant point process to
independent bin counts holds for \emph{any} positive rate
parameterization; the log link is what makes the reduced problem a
canonical-link GLM with the concavity of \cref{thm:ch04:concavity}.
Second, exactness requires \cref{asm:ch04:pwconst}. If covariates vary
within a bin and the fit uses a bin-representative value $\bar\xcov_c$, a
Jensen gap opens --- the true cell mean $\int e^{\bcoef\T\xcov(u)}\dd u$
exceeds $w_c e^{\bcoef\T\bar\xcov_c}$ whenever $\xcov$ genuinely varies
--- and the GLM becomes an approximation. At weekly bins with monthly
macro series and calendar terms the gap is identically zero; the exactness
is a property of our design, not a general fact about binning.

\section{Right censoring: the complete treatment at this rung}\label{ch04:sec:censoring}

What must Poisson calibration do about right censoring? The answer,
proved here in
both of the bookkeeping styles used by the codebase, is: \emph{nothing
beyond keeping the censored exposure}. There is no correction term, no
weighting, no imputation; the survival factor that a censored spell must
contribute is already computed by the compensator integral, provided open
spells are kept and exposure is integrated to the end of observation.

\begin{proposition}[Intensity view: censoring is already in the likelihood]\label{prop:ch04:censorexact}
Let unit $e$ be observed on $[\entry_e, C_e]$ with $C_e = \min(\Tend,
U_e)$, where $\Tend$ is the administrative window end and $U_e$ a recorded
exit time, and let censoring be independent in the sense of
\cref{ch:foundations} (the intensity of $N_e$ with respect to the
filtration enlarged by censoring information is unchanged;
\citealp[III.2]{abgk1993}). Administrative censoring at the deterministic
$\Tend$ is independent by construction; removal at a \emph{recorded} exit
is valid for the cause-specific funding intensity with no independence
assumption at all (\cref{ch:data}, defect (C3)(iii);
\citealp{putter2007}). Then the observed-data log-likelihood is exactly
\cref{eq:ch04:loglik} with $\atrisk_e(u) = \ind{\entry_e \le u \le C_e}$;
in particular the event-free portion $(t_{\mathrm{last}}, C_e]$ of an open
spell contributes
\[
-\int_{t_{\mathrm{last}}}^{C_e} \lam_e(u) \dd u
\]
to $\loglik$ --- the log of its survival factor --- and nothing else. In
the binned form of \cref{thm:ch04:glmequiv}, the entire treatment is the
exposure $w_c = \int_{\text{cell}} \ind{\entry_e \le u \le C_e} \dd u$:
keeping open spells and integrating exposure to $\min(\Tend, U_e)$ is the
complete and exact treatment of (C1).
\end{proposition}

\begin{proof}
Under independent censoring the observed counting process $\int
\ind{u \le C_e} \dd N_e(u)$ has intensity $\atrisk_e(t)\lam_e(t)$ with
respect to the enlarged filtration (\cref{ch:foundations};
\citealp[III.2]{abgk1993}), so the Jacod likelihood of the observed data
is \cref{eq:ch02:jacod} with this intensity --- which is
\cref{eq:ch04:loglik} restricted to $[\entry_e, C_e]$. Splitting the
compensator integral at the observed event times isolates the term over
$(t_{\mathrm{last}}, C_e]$, an interval containing no event: its entire
contribution is $-\int \lam_e$, with no $\log\lam$ term to drop or
invent. The binned statement follows from \cref{thm:ch04:glmequiv}, whose
exposure $w_c$ integrates precisely this at-risk indicator. The
proposition fails exactly when censoring is \emph{informative} ---
observation ends for reasons correlated with the intensity given
covariates. That is not (C1) but (C3), the unlabeled-exit defect: an
unrecorded death is a censoring event we never see and therefore cannot
put into $C_e$; \cref{ch04:sec:defects} states the resulting bias.
\end{proof}

\begin{proposition}[Waiting-time view: the same computation]\label{prop:ch04:renewal}
Under \cref{asm:ch04:condpoisson}, for a spell starting at $a$ (an entry
or an observed event) the waiting time to the next event has survival
function
\begin{equation}\label{eq:ch04:survival}
S(c) \;=\; \Prob\big(N(a, c] = 0\big)
 \;=\; \exp\!\Big(-\int_a^{c} \lam(u)\dd u\Big),
 \qquad c \ge a,
\end{equation}
and density $f(c) = \lam(c)\, S(c)$. The likelihood assembled
spell-by-spell --- density factors for completed spells, survival factors
$S(C)$ for spells censored at $C$ --- equals the Jacod likelihood of
\cref{prop:ch04:censorexact} exactly:
\[
\prod_{m=1}^{n} \lam(t_m)\,
 e^{-\int_{t_{m-1}}^{t_m} \lam(u)\dd u}
 \;\times\;
 e^{-\int_{t_n}^{C} \lam(u)\dd u}
 \;=\;
 \Big(\prod_{m=1}^{n} \lam(t_m)\Big)\,
 e^{-\int_{t_0}^{C} \lam(u)\dd u},
\qquad t_0 = \entry .
\]
\end{proposition}

\begin{proof}
\Cref{eq:ch04:survival} is the Poisson count law of
\cref{asm:ch04:condpoisson} evaluated at zero; differentiating $1 - S(c)$
gives $f(c) = \lam(c) S(c)$. Taking logs of the displayed product, the
survival exponents telescope over the partition $t_0 < t_1 < \cdots < t_n
< C$ into the single compensator integral $-\int_{t_0}^{C}\lam$, and the
event factors contribute $\sum_m \log\lam(t_m)$: this is
\cref{eq:ch04:loglik} for the unit. The two views are therefore not two
treatments to be reconciled but one number computed with two
bookkeepings. Practical corollary: a spell-table implementation (one row
per spell with a censoring flag, survival factor for flagged rows) and an
exposure-table implementation (one row per unit-week cell with exposure
$w_c$) must return identical log-likelihoods on the same data --- a cheap
and mandatory cross-implementation test. In the homogeneous case
\cref{eq:ch04:survival} is the exponential waiting time $e^{-\mub(c-a)}$,
and the renewal view is classical survival analysis of exponential
spells; for recurrent events the counting-process formulation is the one
that generalizes \citep{cook2007}.
\end{proof}

\begin{corollary}[Interval censoring is free at this rung]\label{cor:ch04:intervalexact}
Suppose deal dates are trusted only at resolution $\delta$ (the week): the
bin of each event is correct, its position within the bin is not. Under
the assumptions of \cref{thm:ch04:glmequiv} with bin width $\delta$, the
binned Poisson likelihood \cref{eq:ch04:glm} is the \emph{exact}
likelihood of the trusted data --- not an approximation --- because by
part (iii) the within-bin positions carry no information about $\bcoef$:
the corrupted coordinates are exactly the ancillary ones. Boundary cells
truncated by entry or censoring retain exactness through their partial
exposure $w_c$. None of this survives history dependence: for a Hawkes
process the intensity after $t$ depends on the within-bin positions of
earlier events, binning destroys likelihood information, and the
interval-censored machinery of \cref{ch:censoring}
\citep{rizoiu2022} becomes necessary. The contrast is the precise sense
in which weekly aggregation is free at rung 1 and costly afterwards ---
and it is why (C4) appears as ``robust'' in this chapter's defect table
and as a first-class problem in \cref{ch:hawkes,ch:censoring}.
\end{corollary}

\begin{warning}[Dropping open spells is right truncation, not censoring]\label{warn:ch04:truncation}
The exactness of
\cref{prop:ch04:censorexact,prop:ch04:renewal} holds only if censored
spells are \emph{kept}. Discarding them --- fitting only spells that end
in an observed event --- conditions on the event landing in-window, which
oversamples short gaps and biases every fitted rate upward; the
observation scheme becomes right truncation, a different and harder
problem \citep{kalbfleisch2002}. This restates
\cref{warn:ch02:openspells} at the point of use. For the weekly sector
GLM the trap is nearly impossible to fall into (every sector-week cell
exists whether or not it contains events --- the zero counts \emph{are}
the censored exposure), which is itself an argument for the cell
representation; for spell-table firm-level code the trap is one careless
\code{dropna} away, and the open-spell retention audit of
\modrepo{data/audits.py} exists to catch it.
\end{warning}

\section{Overdispersion: from Poisson to negative binomial}\label{ch04:sec:overdisp}

\Cref{asm:ch04:condpoisson} forces $\Var(y_c) = \E(y_c)$ cell by cell.
We expect funding counts to violate this: latent heterogeneity (unmodeled
sector-week quality, vendor coverage drift) and genuine history dependence
both inflate variance beyond the mean. The synthetic studies of
\cref{ch:results-synth} will measure the dispersion this produces under
controlled misspecification. Two upgrades
present themselves, one inferential and one distributional, and the
design needs both for different purposes.

The inferential upgrade is \emph{quasi-Poisson}: keep the mean model and
the estimating equation \cref{eq:ch04:score}, posit $\Var(y_c) = \phi\,
\mu_c$, estimate $\hat\phi$ by the Pearson statistic
$\smash{\hat\phi = \tfrac{1}{n-d}\sum_c (y_c - \hat\mu_c)^2/\hat\mu_c}$,
and scale standard errors by $\smash{\hat\phi^{1/2}}$ (or use the
sandwich variance, which is robust to variance misspecification beyond
the proportional form). Because the score \cref{eq:ch04:score} is
unbiased whenever the \emph{mean} is correct, $\hat\bcoef$ keeps its
consistency under overdispersion; only its reported precision was
wrong. Quasi-Poisson repairs the standard errors at zero modeling cost, but
it supplies no predictive distribution --- and the evaluation protocols
of \cref{ch:gof} score predictive distributions (PIT, proper scoring
rules), not point forecasts. Hence the distributional upgrade.

\begin{proposition}[Gamma frailty makes a negative binomial]\label{prop:ch04:negbin}
Let $y \st Z \sim \mathrm{Poisson}(Z\mu)$ with frailty $Z \sim
\mathrm{Gamma}(\alpha, \alpha)$ (shape and rate $\alpha$, so $\E Z = 1$,
$\Var Z = 1/\alpha$). Then $y$ is negative binomial:
\begin{equation}\label{eq:ch04:nbpmf}
\Prob(y = k)
 \;=\;
 \frac{\Gamma(k + \alpha)}{\Gamma(\alpha)\, k!}
 \left(\frac{\alpha}{\alpha + \mu}\right)^{\!\alpha}
 \left(\frac{\mu}{\alpha + \mu}\right)^{\!k},
 \qquad k = 0, 1, 2, \ldots,
\end{equation}
with $\E y = \mu$ and $\Var y = \mu\,(1 + \mu/\alpha)$. As $\alpha \to
\infty$ the law converges to $\mathrm{Poisson}(\mu)$. Moreover, with
$\mu_c = w_c e^{\bcoef\T\xcov_c}$ and $\alpha$ fixed, the negative
binomial log-likelihood remains concave in $\bcoef$.
\end{proposition}

\begin{proof}
Integrating the Poisson pmf against the Gamma density,
\[
\Prob(y = k)
 = \int_0^\infty \frac{e^{-z\mu}(z\mu)^k}{k!}\,
 \frac{\alpha^\alpha z^{\alpha-1} e^{-\alpha z}}{\Gamma(\alpha)} \dd z
 = \frac{\mu^k \alpha^\alpha}{k!\,\Gamma(\alpha)}
 \int_0^\infty z^{k+\alpha-1} e^{-(\mu+\alpha)z} \dd z
 = \frac{\mu^k \alpha^\alpha}{k!\,\Gamma(\alpha)}\,
 \frac{\Gamma(k+\alpha)}{(\mu+\alpha)^{k+\alpha}},
\]
which rearranges to \cref{eq:ch04:nbpmf}. The moments follow from the
tower rule: $\E y = \E[Z\mu] = \mu$ and $\Var y = \E[Z\mu] +
\Var(Z\mu) = \mu + \mu^2/\alpha$. For $\alpha \to \infty$, $Z \to 1$ in
probability and the mixture collapses to the Poisson. For concavity,
write $\eta_c = \log w_c + \bcoef\T\xcov_c$; the $\bcoef$-dependent part
of the log-likelihood per cell is $y_c\,\eta_c - (y_c +
\alpha)\log(\alpha + e^{\eta_c})$. The first term is linear, and
$\log(\alpha + e^{\eta}) = \logsumexp(\log\alpha, \eta)$ is convex in
$\eta$, hence $-(y_c+\alpha)\log(\alpha + e^{\eta_c})$ is concave in the
linear function $\eta_c$ of $\bcoef$; sums of concave functions are
concave.
\end{proof}

The mixture derivation is not a convenience trick; it is the model's
meaning. $Z$ is a multiplicative \emph{frailty} --- the latent common
factor of \citet{vaupel1979}, developed for multivariate event data by
\citet{hougaard2000} --- attached here to each sector-week cell to absorb
heterogeneity the covariates miss. That reading also sets the limits of
the upgrade. Independent frailties per cell fix the
\emph{marginal} variance but assert no dynamics; if the missing structure
is autocorrelated (a persistent latent factor, or genuine
event-to-event excitation), Pearson residuals will still be serially
correlated after the NB upgrade, and that residual autocorrelation ---
not overdispersion itself --- is the escalation signal toward
\cref{ch:hawkes}. The distinction matters because overdispersion is
produced by frailty \emph{and} by contagion alike; a dispersion
statistic can demand escalation but can never certify excitation. If
clustered variance is read as contagion, branching estimates inflate;
the frailty experiments of \cref{ch:results-synth} are designed to
measure that inflation. The frailty-vs-contagion
identification problem is treated properly in \cref{ch:hawkes}.

\begin{decision}{D4.2 --- negative binomial is the default count noise}
The default observation model for sector-week counts is the negative
binomial \cref{eq:ch04:nbpmf} with mean $w_c e^{\bcoef\T\xcov_c}$ and a
fitted dispersion $\alpha$ (per sector, unless pooling is justified by
the diagnostics), replacing the pure Poisson wherever a predictive count
distribution is consumed --- held-out NLL, PIT calibration, decision
outputs. Rationale: overdispersion is expected from heterogeneity and
history dependence (to be measured in \cref{ch:results-synth}); an exact
frailty interpretation
(\cref{prop:ch04:negbin}) that absorbs heterogeneity without asserting
dynamics; preserved mean structure and $\bcoef$ interpretation;
preserved concavity in $\bcoef$ at fixed $\alpha$. Quasi-Poisson/sandwich
standard errors remain the reporting standard for coefficient inference
in plain-Poisson fits. Dispersion diagnostics --- $\hat\phi$, fitted
$1/\alpha$, and randomized PIT --- are computed with every count fit
under the protocols of \cref{ch:gof}. Reversal triggers: if real-data
fits drive $1/\alpha \to 0$ after the covariate baseline (equidispersion),
the upgrade is dropped as needless; if NB residuals stay serially
correlated, the deficiency is dynamic and the remedy is history terms
(\cref{ch04:sec:history}, \cref{ch:hawkes}), not a fatter marginal tail.
\end{decision}

\section{History as covariates: the self-modulated GLM, previewed}\label{ch04:sec:history}

One step remains inside the GLM family, and the repository's sector layer
already uses it. Add the recent past to the design: lagged own- and
cross-sector counts (or exponentially filtered versions of them) enter
the linear predictor,
\begin{equation}\label{eq:ch04:autoregression}
\log \Lam_{s,t}
 \;=\;
 a_s + \bcoef\T \xcov_t
 + \sum_{r \in \sectors} \sum_{\ell=1}^{L} b_{sr\ell}\, y_{r, t-\ell},
\end{equation}
which is a \emph{log-linear Poisson autoregression} --- the count time
series family of \citet{fokianos2009}, log-linear variant
\citet{fokianos2011}. Everything mechanical from this chapter survives:
the parameters sit in a linear predictor, so the likelihood is concave
(\cref{thm:ch04:concavity} verbatim), the solver guarantees of
\cref{cor:ch04:global} hold, and the exp link grants \emph{sign freedom}
--- negative $b$ coefficients express inhibition, one of the three legal
routes to it catalogued in \cref{ch:hawkes}. This is rung 3 of the
ladder. Fixed lags (or fixed-decay filters)
play the role of a fixed kernel: they capture the predictive content of
excitation without estimating the timescale, the one parameter
that weekly data cannot identify (the $\branch$--$\decay$ ridge;
\cref{ch:gof}).

Two prices are paid, and both are conceptual rather than computational.
First, lagged counts are \emph{internal} covariates
(\cref{asm:ch04:external} is voided): the model is a legitimate
conditional intensity and one-step-ahead prediction is clean, but
multi-step forecasts and any scenario analysis must simulate the system
jointly --- plugging today's history into a formula for next quarter is
not a probability of anything \citep{kalbfleisch2002}. Second, a
stability condition that must be stated explicitly:

\begin{warning}[Stability of log-link feedback is an operating-point condition]\label{warn:ch04:loglinkstability}
For the \emph{additive} excitation model of \cref{ch:hawkes}, with
nonnegative lag matrices summing to $\Gmat$, subcriticality is the
coefficient condition $\specrad{\Gmat} < 1$, enforced in
\modrepo{sector\_stability.py} through the row-sum bound $\rhomax = 0.95$.
That certificate does \emph{not} transfer to
\cref{eq:ch04:autoregression}. Under the log link one past event
multiplies next week's rate by $e^{b}$, so the marginal effect of an
event on the conditional mean is $b\,\Lam$ --- coefficient \emph{times
operating level}. A feedback matrix that is harmless at $0.3$ events per
sector-week can be explosive at $3$; no inequality on the $b$'s alone,
spectral or otherwise, certifies stability. The applicable theory is the
contraction/ergodicity analysis of \citet{fokianos2011}, whose conditions
constrain the feedback of \emph{bounded or slowly varying} transforms of
the past; raw counts in the exponent sit outside the Lipschitz structure
that the theory requires. We expect simulated log-link feedback models to
produce metastable excursions even when the nominal spectral radius is
subcritical; the stress tests of \cref{ch:results-synth} probe exactly
this. Design rules that follow:
history enters exponents only through bounded transforms (EWMA levels,
$\log(1+\cdot)$ counts); any simulation or multi-step use of a fitted
self-modulated GLM must verify stability \emph{at the operating point};
and the branching reading of excitation is reserved for the additive
parameterization. The full treatment --- both parameterizations, their
stability certificates, and what each can and cannot claim about
contagion --- is \cref{ch:hawkes}.
\end{warning}

\section{Defect declaration}\label{ch04:sec:defects}

Per Design decision D2.1 (\cref{ch:data}), the Poisson models of this
chapter declare their position against the defect ledger of
\cref{tab:ch02:defects}. \Cref{tab:ch04:defects} summarizes; the signs
are the chapter's own results.

\begin{table}[htbp]
\centering
\small
\begin{tabular}{p{0.06\textwidth} p{0.26\textwidth} p{0.36\textwidth} p{0.18\textwidth}}
\toprule
\textbf{Defect} & \textbf{Status at this rung} & \textbf{Mechanism and
sign} & \textbf{Treatment} \\
\midrule
(C1) & corrected \emph{exactly} & censored exposure contributes its
survival factor through the compensator; no bias
(\cref{prop:ch04:censorexact}) & this chapter \\
\addlinespace
(C2) & corrected mechanically & risk sets start at $\entry_i$; exposure
$w_c$ integrates from entry. Residual selection (``firms that have
raised'') is an interpretation limit, not a likelihood error &
\cref{ch:universe} \\
\addlinespace
(C3) & \emph{vulnerable} (firm level) & unrecorded exits inflate exposure
$w_c$ with dead time; fitted hazards \emph{deflated} (downward sign),
worst for long-quiet firms. Sector level immune: sectors do not die &
exposure windows, \cref{ch:buckets} \\
\addlinespace
(C4) & robust at weekly resolution & binned likelihood is exact; within-bin
positions are ancillary (\cref{cor:ch04:intervalexact}); day-level jitter
within weeks is immaterial & \cref{ch:censoring} \\
\addlinespace
(C5) & \emph{vulnerable} & stale firm covariates attenuate their fitted
effects toward zero; macro covariates unaffected (published as-of, never
stale). Bites the event-time firm path, not the sector GLM &
\cref{ch:stale} \\
\addlinespace
(C6) & sector immune / firm vulnerable & sector counts are complete
(out-of-universe deals still carry sectors); firm-attributable intensity
misses competitors and thins counts & \cref{ch:universe} \\
\bottomrule
\end{tabular}
\caption{Defect declaration for the Poisson rung (D2.1 compliance).
``Immune''/``robust'' claims hold under this chapter's stated
assumptions; signed vulnerabilities are accepted and carried to the
chapters that own the remedies.}
\label{tab:ch04:defects}
\end{table}

The asymmetry in \cref{tab:ch04:defects} between sector and firm rows is
the quantitative reason the ground process is calibrated at sector
resolution first: at that resolution the Poisson rung is exactly right
under (C1), (C2), (C4) and immune to (C3), (C6), so its held-out
likelihood is a trustworthy floor. Firm-level Poisson calibration through
the event-time path inherits the (C3) deflation and (C5) attenuation with
the stated signs, and every later firm-level model must beat it
\emph{while carrying the same declared exposures}.

\begin{codehooks}
\emph{Existing.} \modrepo{sector\_stability.py}
(\code{fit\_sector\_count\_model}): with \code{n\_lags=0} the additive
excitation vanishes and the fit is exactly the weekly NHPP GLM of
\cref{thm:ch04:glmequiv,cor:ch04:repoglm}; its ridge \code{l2\_beta} and
baseline support-envelope winsorization implement the remedies of
\cref{warn:ch04:separation}; \code{sector\_rate\_at} is the one-step
predictor. \code{synthetic/fast\_fit.py::fit\_sector\_glm\_fast(n\_lags=0)}
is the vectorized synthetic-world twin for the decomposition benchmarks
planned in \cref{ch:results-synth}. \modrepo{eventtime/estimate.py}
(\code{fit\_multivariate\_with\_covariates}) is the day-resolution NHPP
path when excitation is disabled, with
\modrepo{eventtime/covariates.py}
(\code{PiecewiseConstantCovariate.integrate\_exp\_gamma}) computing the
compensator $\int e^{\bcoef\T\xcov}$ in closed form cell by cell ---
the continuous-time twin of the GLM offset. \emph{To be written.}
\modrepo{sector\_negbin.py}: the D4.2 negative-binomial layer ---
NB likelihood \cref{eq:ch04:nbpmf} with log-mean $\log w_c +
\bcoef\T\xcov_c$, per-sector $\alpha$, quasi-Poisson $\hat\phi$ and PIT
diagnostics emitted in the \cref{ch:gof} harness format.
\modrepo{data/audits.py} (exposure integrity, per \cref{ch:data}):
open-spell retention (\cref{warn:ch04:truncation}), entry-date
consistency of $w_c$ integration limits, spell-table vs.\ exposure-table
log-likelihood parity (\cref{prop:ch04:renewal}), and train/test
covariate-support overlap reporting (\cref{warn:ch04:separation}).
\end{codehooks}
